\section{Quantum Dependence}\label{sec:dependence}

This section defines what a slice of a quantum circuit must preserve, what an operation's dependence
on a qubit is, and the dependence rules we compare. The definitions are independent of any slicer, and
the evaluation in Section~\ref{sec:results} measures each rule against them.

\subsection{Circuits and their semantics}
A circuit $P = o_0 \dots o_{L-1}$ is a sequence of operations over qubits $Q = \{0, \dots, n-1\}$ and
classical bits $C$. Each operation $o$ acts on a set of \emph{locations} $\mathit{loc}(o) \subseteq
Q \cup C$. A gate acts on its qubits, a measurement on its qubit and the bit it writes, a reset on its
qubit, and a classically conditioned gate on its qubits and the bits its condition reads. For an input
state $\rho$ of the qubits, with every bit initially $0$, $\llbracket P \rrbracket(\rho)$ denotes the
unconditional final state, that is, the joint state of qubits and bits averaged over measurement
outcomes. For $W \subseteq Q$, $\rho_W = \operatorname{Tr}_{Q \setminus W}(\rho)$ is the reduced
state of $W$. For a set of positions $S$, the restriction $P|_S$ keeps the operations at the positions
in $S$, in their original order. Every slicer in this paper returns such a set, and every slice is
judged by running $P|_S$.

\subsection{Criteria, slices, and dependence}

\begin{definition}[Slicing criterion]\label{def:criterion}
A slicing criterion is a nonempty set $W \subseteq Q$ of qubits at the end of the circuit.
\end{definition}

The classical criterion $(s, V)$ names a statement $s$ and variables $V$ used at $s$, and Yu and Zhao
transfer it unchanged to quantum programs~\cite{yu2025qpdg}. A single qubit has no value of its own,
however: its state can be determined only together with the qubits it is correlated with. We therefore
let the criterion be a set of qubits and preserve their \emph{joint} state. $W = \{q\}$ is the
single-qubit criterion.

\begin{definition}[State preservation and soundness]\label{def:sound}
A set of positions $S$ \emph{preserves} $W$ at input $\rho$ if
$\operatorname{Tr}_{Q \setminus W}\llbracket P|_S \rrbracket(\rho) =
\operatorname{Tr}_{Q \setminus W}\llbracket P \rrbracket(\rho)$. $S$ is a \emph{sound} slice for $W$ if
it preserves $W$ at every input.
\end{definition}

Comparing reduced states is basis-independent. It detects phase errors that a measured distribution in
one basis misses. In the experiments, equality means Uhlmann fidelity at least $1 - 10^{-6}$.

Soundness says that the slice reproduces what the criterion qubits end up in. A developer localizing a
fault needs a second guarantee, that the slice contains every place where a fault could change that
outcome.

\begin{definition}[Quantum dependence]\label{def:dep}
Operation $o_k$ is a \emph{dependence} of $W$ in $P$ if there exist an input $\rho$ and a unitary $U$ on
the qubits of $o_k$ such that inserting $U$ immediately after $o_k$ changes
$\operatorname{Tr}_{Q \setminus W}\llbracket P \rrbracket(\rho)$. For a measurement, inverting its
recorded outcome for the conditions that read it also counts as such a change.
$\mathit{Dep}_W(P)$ is the set of dependences of $W$. A set $S$ is \emph{dependence-complete} for $W$ if
$\mathit{Dep}_W(P) \subseteq S$.
\end{definition}

The two properties are distinct. Deleting a single unitary $o_k \notin \mathit{Dep}_W(P)$ leaves
$\rho_W$ unchanged (take $U = o_k^{\dagger}$), but deleting several such operations can change it. A
sound slice, conversely, need not contain every dependence. The ground truth of
Section~\ref{sec:design} approximates the existential quantifiers of Definition~\ref{def:dep} by sampled
inputs and perturbations, so the measured $\mathit{Dep}_W(P)$ is a subset of the true one.

\subsection{Dependence rules}

A \emph{dependence rule} $R$ assigns each operation the locations it reads, $r_R(o)$, and the
locations it writes, $w_R(o)$. Operation $o_j$ depends on $o_i$ under $R$ if $i < j$ and some location
in $w_R(o_i) \cap r_R(o_j)$ is not written by any operation between them. The backward slice
$B_R(W)$ is the set of operations from which the last writer of some $q \in W$ is reachable. A rule can
account for three kinds of dependence. \emph{Data} dependence follows a qubit from the operation that
last changed it. \emph{Control} dependence follows a classical bit from the measurement that wrote it to
the operations its value conditions. \emph{Quantum} dependence lets a multi-qubit gate change every qubit
it acts on, including its controls. We compare three rules that differ in which of these they include,
and two refinements of the most complete one.

\emph{Classical PDG} (CPDG) applies the classical program dependence graph~\cite{ferrante1987pdg} to the
circuit, with data and control dependence. Its data dependence reads a gate as a classical assignment. A
controlled gate reads its controls and targets and writes only its targets, and any other gate reads and
writes all its qubits. Its control dependence follows the classical bits. A measurement writes its bit,
and a conditioned gate reads the bits its condition tests. We also report \emph{def-use} (DU), the same
data dependence without classical bits, as the minimal classical port.

\emph{Wire} (W) adds quantum dependence to the classical PDG. Every operation reads and writes all its
locations, $r_W(o) = w_W(o) = \mathit{loc}(o)$, so a controlled gate also writes its controls. $B_W$ is the
backward light cone of $W$, extended along classical bits from measurements to the gates they condition.

\emph{Entanglement closure} (QPDG) is our documented reading of Yu and Zhao's
graph~\cite{yu2025qpdg}, whose dependences are given informally. It adds to the wire edges an edge from
every operation to every later one when both act on qubits of the same component of the coupling
graph, which links any two qubits joined by a multi-qubit gate anywhere in the circuit.

\emph{Commutation} keeps the operations of $B_W(W)$ that do not commute with the observables
propagated backward from $W$, the union over the Pauli observables $X$, $Y$ and $Z$ on each criterion
qubit, as implemented by Qiskit's \texttt{LightCone} transpiler pass~\cite{javadi2024qiskit}.

\emph{State-based pruning} (QSlice) starts from $B_W(W)$ at a given input, by default $\ket{0\dots0}$.
It removes an operation that does not act on the criterion when the mutual information between the
criterion and the operation's qubits is below $10^{-9}$, and when the removal still preserves $W$ at
that input.

\subsection{Properties}

\begin{theorem}[The wire rule is sound]\label{thm:sound}
For every circuit $P$ and criterion $W$, $B_W(W)$ is a sound slice for $W$.
\end{theorem}

\begin{IEEEproof}
Process $P$ backward with a set of needed locations $N$, initially $W$. Operation $o_k$ joins the slice,
and $N$ grows by $\mathit{loc}(o_k)$, exactly when $\mathit{loc}(o_k) \cap N \neq \emptyset$. This
computes $B_W(W)$. An operation outside the slice acts only on locations outside the current $N$. Such a
trace-preserving operation leaves the reduced state on $N$ unchanged (no-signaling), so removing it
preserves the reduced state of the needed locations at that point, for every input. Induction over
$P$ gives the claim for $W$. The theorem is mechanized in Isabelle/HOL for an abstract model of local
operations whose assumptions are shown satisfiable.
\end{IEEEproof}

\begin{proposition}[The wire rule is dependence-complete]\label{prop:complete}
$\mathit{Dep}_W(P) \subseteq B_W(W)$.
\end{proposition}

\begin{IEEEproof}
Let $o_k \notin B_W(W)$, and let $P'$ insert $U$ after $o_k$. If $U$ were in $B_W(W)$ for $P'$,
then $o_k$, which precedes $U$ on the same locations, would be in $B_W(W)$ for $P$. So $P'$ and $P$
have the same wire slice, and by Theorem~\ref{thm:sound} both preserve $W$. Hence
$\operatorname{Tr}_{Q \setminus W}\llbracket P' \rrbracket(\rho) =
\operatorname{Tr}_{Q \setminus W}\llbracket P|_{B_W(W)} \rrbracket(\rho) =
\operatorname{Tr}_{Q \setminus W}\llbracket P \rrbracket(\rho)$. The argument for an inverted
measurement outcome is the same.
\end{IEEEproof}

The other rules fail one property or the other, as the two programs of Fig.~\ref{fig:dependence} show.
Each node of their graphs is an operation, labelled with the source line that created it, and each edge
is a dependence of the wire rule, marked as data, control or quantum.

Control dependence is needed whenever measurement outcomes steer later gates. In the teleportation
program of Fig.~\ref{fig:dependence}a, the corrections on $q_2$ run only if the measurements of $q_0$
and $q_1$ return $1$. Def-use ignores the classical bits. Its slice of $q_2$ keeps \numTelDuOps{} of the
\numTelOps{} operations, and its fidelity is \numTelDuFidelity{}. The classical PDG follows the control
edges from the measurements to the conditioned gates, and its slice equals the wire slice.

\begin{proposition}[The classical PDG is unsound]\label{prop:du}
There are circuits for which $B_{CPDG}(W)$, and hence $B_{DU}(W)$, does not preserve $W$ even at
$\ket{0\dots0}$.
\end{proposition}

In the ansatz of Fig.~\ref{fig:dependence}c, each \texttt{cx(qubit, qubit + 1)} acts on its control as
well as its target. After the first layer, the state of $q_0$ depends on the rotations of $q_1$ and $q_2$
through these gates. The classical PDG treats a control as read only, so its slice of $q_0$ keeps only
the \numAnsCpdgOps{} rotations on $q_0$ itself, the bold nodes of Fig.~\ref{fig:dependence}d, with
fidelity \numAnsCpdgFidelity{}. The edges it lacks are exactly the quantum edges of the graph. Control
dependence cannot repair this, because the missing dependence is quantum. The wire slice keeps the
\numAnsWireOps{} operations that are dependences of $q_0$, and no other.

\begin{proposition}[Entanglement is neither necessary nor sufficient for dependence]\label{prop:ent}
(i) An operation can be a dependence of $q$ although $q$ is never entangled with its qubits.
(ii) An operation can act on a qubit entangled with $q$ and not be a dependence of $q$.
\end{proposition}

For (i), $q_0$ in the buggy program of Fig.~\ref{fig:motivating} is in a pure state after every
operation, so it is never entangled with anything. Yet \texttt{h(1)} and the CNOT are dependences of
$q_0$, since the fix of line~6 changes the state of $q_0$. For (ii), in
$\mathrm{H}(q_0)\,\mathrm{CX}(q_0, q_1)\,\mathrm{X}(q_1)$ the two qubits are maximally entangled, but
$\mathrm{X}(q_1)$ is not a dependence of $q_0$. No operation on other qubits can change $\rho_{q_0}$
(no-signaling). Entanglement closure therefore adds no dependence that the wire rule misses. It only adds
operations that are not dependences. In Fig.~\ref{fig:dependence}d it adds the \numAnsYzExtra{} rotations
of $q_2$ in the second layer, which run after the last gate that can carry information from $q_2$ to $q_0$.

\begin{proposition}[The wire rule distributes over criteria]\label{prop:joint}
$B_W(W_1 \cup W_2) = B_W(W_1) \cup B_W(W_2)$.
\end{proposition}

\begin{IEEEproof}
By induction over the backward pass of Theorem~\ref{thm:sound}, the needed set for $W_1 \cup W_2$ is
the union of the needed sets for $W_1$ and $W_2$ at every position.
\end{IEEEproof}

Under the wire rule, a joint criterion is therefore the union of single-qubit slices. The criterion
matters for pruning. State-based pruning for $\{q\}$ preserves $\rho_q$ but not the correlations of $q$
with the other qubits, so it can delete an operation that changes $\rho_W$ for a larger $W$. It is also
specific to its input, because the mutual information it tests and the reduced state it checks are
those of one input. Section~\ref{sec:results} measures both effects.

\paragraph*{Superposition}
None of the rules depends on amplitudes. A dependence graph has one node per operation and edges
between operations on shared locations, so superposition cannot enlarge it~\cite{nielsen2010quantum}.
Superposition matters only for state-based pruning, whose checks simulate the circuit and cost time
exponential in its width (Section~\ref{sec:results}).
